\subsection{Graph cofibers over corners of order at most six}

\begin{theorem}[Proved: the bounded small-corner graph obstruction]
\label{CP2:A-baseline-six-graph}
Let a finite monoid $M$ of order at most six act on the
right of a nonempty finite graph $\Gamma$.  Require each
actor to take a vertex to a vertex and each parametrized
edge to an edge, preserving or reversing its unit parameter,
or to a vertex.  If
\[
 X=\operatorname{hocofib}(\Gamma_{hM}\longrightarrow BM)
\]
has bounded rational homology, every rational degree-two
product on $X$ is zero.  The left-action statement is dual.
\end{theorem}
\begin{proof}
If $H^2(BM;\mathbb Q)=0$, each second class is relative
to $BM$ and lifts from the suspension quotient of the
cofiber.  Its products vanish.  Otherwise the minimal
ideal must have at least two rows and two columns;
the small cardinality forces a $2\times2$ rectangular
band $K$, with trivial maximal group and at most two
outside elements.  Put $R=\mathbb QM$ and
$A=R/\mathbb QK$.

Choose $e=x_{00}$ and $e_1=x_{01}$.  Their image graphs
$\Gamma_0=\Gamma e$ and $\Gamma_1=\Gamma e_1$ are
retract subgraphs.  An image under a band element depends
only on its column, and right multiplication by $e_1$
and $e$ gives mutually inverse graph isomorphisms between
these two images.  If they share a vertex $v$, then
$ve=ve_1=v$ and
$vx_{ij}=vex_{ij}=vx_{0j}=v$.  For every $m\in M$,
$vm=v(em)=v$ because $em\in K$.  This common fixed
vertex gives an action-bar section of
$\Gamma_{hM}\to BM$.  The cofiber sequence makes every
positive cohomology class relative to $BM$, and the
suspension quotient again kills its products.

Suppose instead that the two images have disjoint vertex
sets, hence disjoint edge sets.  Use the augmented graph
complex of right $R$-modules
\[
 C=[C_1(\Gamma;\mathbb Q)\longrightarrow
                  \ker(C_0(\Gamma;\mathbb Q)\to\mathbb Q)]
\]
in degrees one and zero.  The multiplication map
$Ce\otimes_{\mathbb Q}eR\to C$ is injective in both
degrees.  Indeed, $eR$ has the two column basis elements.
In degree one their images have disjoint edge supports;
orientation reversals change only signs.  In degree zero,
$Ce$ consists of the augmentation-zero vertex chains in
$\Gamma_0$, and its two column images have disjoint
vertex supports and each has augmentation zero.  Thus
there is no nonzero relation between the two copies.

Its image is exactly $C\mathbb QK$, using the literal
band factorization $\mathbb QK=Re\otimes_{\mathbb Q}eR$.
Consequently
\[
 D=C/(Ce\otimes_{\mathbb Q}eR)
\]
is an actual complex of right $A$-modules concentrated
in degrees zero and one.  The subcomplex being divided
out is a complex of projective right $R$-modules; after
tensoring with the trivial left module it is just $Ce$,
still in those two degrees.  Therefore
\[
 H_j(C\otimes_R^{\mathbf L}\mathbb Q)
  =H_j(D\otimes_R^{\mathbf L}\mathbb Q),\qquad j\geq3.
\]

Let $V$ be the one-dimensional row augmentation module.
The exact sequence $0\to V\to Re\to\mathbb Q\to0$
and the equality $De=0$ give
$D\otimes_R^{\mathbf L}\mathbb Q
 \simeq(D\otimes_R^{\mathbf L}V)[1]$.
There is no higher change-of-rings correction: the
left-projective resolution
$0\to\mathbb QK\to R\to A\to0$ gives
$U\otimes_R^{\mathbf L}A\simeq U$ for every right
$A$-module $U$, since $Ue=0$.  Applying this naturally
to the two terms of $D$ identifies
$D\otimes_R^{\mathbf L}V$ with
$D\otimes_A^{\mathbf L}V$.
Finally, the reduced graph-cofiber chains are
$(C\otimes_R^{\mathbf L}\mathbb Q)[1]$, by the
augmented graph sequence.  Combining the two shifts
gives the precise formula
\[
 H_n(X;\mathbb Q)
   =H_{n-2}(D\otimes_A^{\mathbf L}V),\qquad n\geq4.
\]
This calculation applies also when $\Gamma$ is disconnected.

The algebra $A$ has dimension one or two.  If it is
semisimple, $V$ is projective and its tensor with $D$
has no homology in degree two, so $H_4(X;\mathbb Q)=0$.
The only other case here is
$A=\mathbb Q[\varepsilon]/(\varepsilon^2)$: with one
outside element $a$ beyond the identity, its square
in the quotient is $0$, $a$, or $1$, the last two
giving semisimple algebras.  In the dual-number case
$V$ is its residue module $k=\mathbb Q$.

Only an elementary two-term recurrence is needed.
Tensor $0\to k\to A\to k\to0$ derivedly with $D$.
The middle tensor is $D$, with zero homology in
degrees at least two.  Its long exact sequence gives
\[
 H_j(D\otimes_A^{\mathbf L}k)
  \cong H_{j-1}(D\otimes_A^{\mathbf L}k),\qquad j\geq3.
\]
Boundedness of $X$, through the displayed localization
formula, makes these groups eventually zero.  The
recurrence forces the degree-two group to be zero,
again giving $H_4(X;\mathbb Q)=0$.  Thus either the
fixed-vertex relative argument or fourth-homology
vanishing annihilates every degree-two product.
\end{proof}

\begin{lemma}[Proved: an equivariant graph for a semisimple pair]
\label{CP2:A-baseline-semisimple-pair-graph}
Let $D$ have two objects with arrows only from the first
to the second and with semisimple rational corner
algebras $\mathbb QN,\mathbb QP$.  For any finite
right $D$-act $F$, its category of elements admits an
actual natural rational homology equivalence to a
finite graph.  Every commuting monoid action induces
strict unit-speed graph maps and makes this comparison
equivariant.  Also $H_j(BD;\mathbb Q)=0$ for $j\geq2$.
\end{lemma}
\begin{proof}
Put $U=F(0)$, $T=F(1)$ and $Z=D(0,1)$, viewed as
a left-$P$, right-$N$ biset.  The graph has vertex
sets $U/N$ and $T/P$ and edge set
$(T\otimes_P Z)/N$.  An edge $[t,z]$ joins $[tz]$
to $[t]$.  Here an orbit set is the quotient generated
by the action relations, including noninvertible actors.
The balanced relations $(tp,z)\sim(t,pz)$ and
$(t,zn)\sim(t,z)$ respect both endpoints.

The functor from the element category sends objects to
these vertex orbits, corner arrows to identities, and a
cross arrow $(t,z)$ to its displayed edge.  The balanced
relations are exactly what composition with corner
arrows requires.  No two cross arrows compose.  Thus
this is an actual functor, with no choice of orbit
representatives.  A commuting action preserves the
relations and sends each graph edge to an edge by
its parameter, so the comparison is equivariant.

Filter element-category bar chains by the number of
strict object changes, which is zero or one.  The
vertex complexes are
$\mathbb QU\otimes_{\mathbb QN}^{\mathbf L}\mathbb Q$
and
$\mathbb QT\otimes_{\mathbb QP}^{\mathbf L}\mathbb Q$.
The one-cross-arrow complex is
$\mathbb QT\otimes_{\mathbb QP}^{\mathbf L}\mathbb QZ
 \otimes_{\mathbb QN}^{\mathbf L}\mathbb Q$.
Semisimplicity puts their homology in degree zero,
where they are exactly the linearizations of the
two vertex orbit sets and the balanced edge set.
Their boundary maps are the actual displayed
endpoints.  The filtration therefore proves that
the functor induces the asserted rational homology
equivalence.  Apply the same construction to the
singleton act to obtain the last assertion about $BD$.
\end{proof}

\begin{corollary}[Proved: a semisimple pair opposite a corner through order six]
\label{CP2:A-baseline-semisimple-pair-six}
Let a three-object category have full linear acyclic
support.  If its first two corner algebras are semisimple
and its last corner has order at most six, bounded
rational homology implies that every rational degree-two
product is zero.  The opposite statement also holds.
In particular the size conditions $a,b\leq2$, $c\leq6$
exclude every projective-plane profile, regardless of
the three cross-hom sizes.
\end{corollary}
\begin{proof}
Use the first two objects as the source block $D$ and
the last corner as $M$.  The two-block model is
$BC\simeq\operatorname{hocolim}(BD\leftarrow E\to BM)$,
where $E$ is the homotopy-orbit space of the incoming
act's element category.  Collapse $BD$ and obtain
$X=\operatorname{hocofib}(E\to BM)$.
Since $H^2(BD;\mathbb Q)=0$, this actual quotient map
induces a surjection $H^2(X;\mathbb Q)\to H^2(BC;\mathbb Q)$.
Since all $H_j(BD;\mathbb Q)$ for $j\geq2$ vanish,
the cofiber sequence also passes bounded rational
homology from $BC$ to $X$.

The preceding lemma provides an equivariant rational
comparison of the element category to a graph $\Gamma$.
Its comparison remains a rational homology equivalence
after taking homotopy orbits: the nonnegative free
monoid bar preserves a coefficient quasi-isomorphism.
It gives an actual rational homology equivalence from
$X$ to $\operatorname{hocofib}(\Gamma_{hM}\to BM)$.
The bounded graph theorem, in its left-action version,
kills all second products there.  Pulling back through
the rational ring isomorphism and then through the
surjection in degree two proves the claim.  Corner
monoids of orders one and two have semisimple rational
algebras, giving the stated profile consequence.
\end{proof}

\begin{corollary}[Proved: the right-zero-three endpoint]
\label{CP2:A-baseline-right-zero-three-six}
Let $R_3$ be the identity-adjoined two-element right-zero
band.  A collage between $R_3$ and any monoid of order
at most six, with the $R_3$ acting on the left of the
connecting biset, has zero rational degree-two products
whenever its rational homology is bounded.  There is
no restriction on the number of connecting arrows.
\end{corollary}
\begin{proof}
The two band actors on the biset are retractions with
common image $U$.  Their action category has the
natural graph model used in the proof of
Lemma~\ref{CP2:A-high-P5-graph}: vertices are $U$,
and an outside state gives the edge between its two
retraction values.  Every commuting map sends such
an edge at unit speed to an edge or to a vertex.
Since $BR_3$ is contractible, the collage is the
graph cofiber of the other monoid's action.  Apply
Theorem~\ref{CP2:A-baseline-six-graph}.
\end{proof}
